\chapter{Graph-Design Screening}
\label{ch:experiments}

This chapter reports the graph-design screens. All are balanced full factorials
on \texttt{bus14} with three seeds per cell, and all ask the same kind of
question: does a specific choice in the actor of Chapter~\ref{ch:methods}
change how well a decentralized MAPPO policy survives on held-out chronics?

\begin{itemize}
  \item \textbf{Screen A} varies the depth and width of the actor encoder,
  which applies to every graph schema: a $3\times4$ grid, 36 runs.
  \emph{Training at the time of writing.}
  \item \textbf{Screen B} varies the message-passing operator and the readout,
  GINE against GAT and mean against maximum pooling: a $2\times2$ factorial,
  9 new runs plus one reused cell. \emph{Training at the time of writing.}
  \item \textbf{Screen C} varies the three structural augmentations of the
  busbar graph -- same-substation edges, substation-summary nodes, and
  virtual-node readout (Sections~\ref{sec:substation-representation}
  and~\ref{sec:virtual-node-readout}): a $2^3$ factorial, 24 runs.
  \emph{Complete.}
  \item \textbf{Screen D} varies the direction of the generator and load
  attachment relations in the disaggregated graph
  (Section~\ref{sec:heterogeneous-graph-representation}): a $3\times3$
  factorial, 27 runs. \emph{Complete.}
\end{itemize}

\noindent Screens A, B, and C all run on the busbar graph and share one reference
configuration, so they can be read as a sequence; only Screen D changes the
schema. The two completed screens account for 51 finished training runs.

\section{How This Chapter Is Organized}
\label{sec:screening-how-to-read}

Each screen is presented as one block containing its design and its outcome
together:

\begin{enumerate}
  \item \textbf{Question.} Which design choice is uncertain?
  \item \textbf{Design.} Which factors are varied, at which levels, and what is
  held fixed.
  \item \textbf{Coverage.} How far the runs actually got, before any endpoint
  number is read.
  \item \textbf{Result.} The measured tables.
  \item \textbf{Reading.} What follows, and what does not.
\end{enumerate}

\section{Shared Protocol}
\label{sec:screening-common-protocol}

Every screen uses the configuration in
Table~\ref{tab:screening-common-settings}. Every setting not named as a
screened factor is identical across all runs, so each screen is a paired
comparison against its own baseline cell.

\begin{table}[H]
  \centering
  \small
  \caption{Settings shared by all screens. Only the graph type and the screened
  factors differ between them.}
  \label{tab:screening-common-settings}
  \setlength{\tabcolsep}{4pt}
  \begin{tabular}{L{0.24\textwidth}L{0.68\textwidth}}
    \toprule
    Component & Fixed setting \\
    \midrule
    Environment & \texttt{bus14}, difficulty 0, topology actions, decentralized agents, normalized observations and rewards, 80/20 train/test chronic split with split seed 0 \\
    Actor & One GNN encoder shared by all agents; LayerNorm; node pre-encoder; two 128-unit action-head layers; one-hop neighbor context \\
    Critic & Centralized MLP with three 256-unit hidden layers \\
    Graph preprocessing & None: neither deterministic physical scaling nor running standardization \\
    Training budget & \(15{,}000{,}000\) environment steps, 72 parallel environments, 576 rollout steps \\
    MAPPO update & 10 epochs, 16 minibatches, actor and critic learning rates \(10^{-4}\) with linear annealing, \(\gamma=0.9\), \(\lambda=0.95\), clip \(0.2\), target KL \(0.02\), entropy coefficient \(0.01\) \\
    Exploration & No initial do-nothing prior and no action-0 logit bonus \\
    Sparse control & Disabled: no intervention gate, no intervention penalty, no evaluation-time rho heuristic \\
    Evaluation & Deterministic evaluation every 82{,}944 environment steps, 10 episodes on each split \\
    Seeds & 0, 1, 2 in every cell \\
    \bottomrule
  \end{tabular}
\end{table}

{\color{blue}\subsection{Reference configuration}
\label{subsec:screening-reference-config}

Screens A, B, and C are all derived, seed by seed, from one reference
configuration: the unaugmented busbar graph with a GINE encoder, two
message-passing layers, 128 hidden and output features, and mean readout. It
appears in the repository as
\texttt{gs\_s2\_bus\_n0\_none\_e0n0v0} and is the \texttt{e0n0v0} cell of
Screen C, where its measured survival is reported
(Section~\ref{sec:screen-bus-structure}). Each of the three screens changes one
group of settings relative to it and leaves everything else untouched, so the
reference cell is the common origin of the whole busbar-graph part of this
chapter. Screen D uses the disaggregated schema and has its own baseline cell.}

\subsection{Endpoint statistics}
\label{subsec:screening-endpoint-statistics}

Survival is the percentage of an episode's maximum length that the agent keeps
the grid alive, reported on the test split. Each run saves the checkpoint with the highest test survival seen during training. That checkpoint is then re-evaluated deterministically over all 201 held-out test chronics.

\subsection{Compute}
\label{subsec:screening-budget-alignment}

The runs were executed on two clusters, JED (CPU) and IZAR (GPU); seeds of the same
cell are spread across both, so the compute backend is not confounded with any
screened factor.

\section{Screen A: Encoder Depth and Width}
\label{sec:screen-depth-width}

\emph{Runs in progress. The design below is final; the coverage, result, and
reading subsections are placeholders to be filled from the completed runs.}

\subsection*{Question}

{\color{blue} The reference configuration uses two message-passing layers and 128 hidden features. Those values were fixed by convention rather than measured. Depth and width control different things: with $K$ layers a node embedding can only see
$K$ hops away, and a local actor graph with one-hop context has a small diameter, so a third layer may add nothing that the readout can use; width sets the capacity of each node embedding and of the graph vector that the action head consumes. Is the default $(2,128)$ actor justified, is one hop enough, and is there a width below which the policy degrades?} 
\\
\\
-> We need to fix the hyperparameter for the message passing depth and the encoder width. This is the most general question in the chapter: depth and width apply to every schema, so the answer conditions how the later screens should be read.

\subsection*{Design}

A complete $3\times4$ grid with three seeds per cell, 36 runs.
Table~\ref{tab:depth-width-design} lists the screened parameters and their
levels. The graph output dimension is matched to the hidden dimension, so a
cell varies one width rather than two. Every other setting is inherited from
the reference configuration of
Section~\ref{subsec:screening-reference-config}, seed by seed.

\begin{table}[H]
  \centering
  \small
  \caption{Screen A. Screened parameters and the values tested. All other
  settings are held at Table~\ref{tab:screening-common-settings}.}
  \label{tab:depth-width-design}
  \setlength{\tabcolsep}{5pt}
  \begin{tabular}{L{0.30\textwidth}L{0.24\textwidth}L{0.30\textwidth}}
    \toprule
    Parameter & Setting & Values tested \\
    \midrule
    Message-passing layers $K$ & \texttt{gnn\_layers} & 1, 2, 3 \\
    Encoder width $d_h$ & \texttt{gnn\_hidden\_dim} & 16, 32, 64, 128 \\
    \bottomrule
  \end{tabular}
\end{table}

\subsection*{Coverage}

\pending{completed steps per run, common budget, and whether any cell lost a
seed. Small-width runs are cheaper per step, so the walltime limit is expected
to bind unevenly across the width axis --- if it does, report the common budget
per width column rather than one budget for the whole grid.}

\subsection*{Result}

\begin{table}[H]
  \centering
  \small
  \caption{Screen A. Test survival by message-passing depth and encoder width,
  averaged over three seeds, with the seed-to-seed standard deviation in
  parentheses. \pending{fill from the completed runs; state which endpoint
  statistic of Section~\ref{subsec:screening-endpoint-statistics} is used.}}
  \label{tab:depth-width-grid}
  \setlength{\tabcolsep}{6pt}
  \begin{tabular}{lcccc}
    \toprule
    & \multicolumn{4}{c}{Encoder width $d_h$} \\
    \cmidrule(lr){2-5}
    Layers $K$ & 16 & 32 & 64 & 128 \\
    \midrule
    1 & -- & -- & -- & -- \\
    2 & -- & -- & -- & -- \\
    3 & -- & -- & -- & -- \\
    \bottomrule
  \end{tabular}
\end{table}

\begin{table}[H]
  \centering
  \small
  \caption{Screen A. Replication check: the $(2,128)$ cell against the
  \texttt{e0n0v0} result of Screen C, which is the same configuration launched
  separately. \pending{fill; a gap here is run-to-run noise, not an effect.}}
  \label{tab:depth-width-replication}
  \setlength{\tabcolsep}{6pt}
  \begin{tabular}{lrrr}
    \toprule
    Launch & Mean (\%) & Std (pp) & Seeds \\
    \midrule
    Screen C \texttt{e0n0v0} & 97.31 & 1.22 & 3 \\
    Screen A $(K,d_h)=(2,128)$ & -- & -- & 3 \\
    \midrule
    Difference & -- & & \\
    \bottomrule
  \end{tabular}
\end{table}

\begin{figure}[H]
  \centering
  \pendingplot{Aggregated test episodic-survival curves, one panel per depth
  $K$, four width curves inside each panel, one standard deviation band across
  the three seeds.}
  \caption{Screen A. Survival curves by depth and width.
  \pending{replace the box with the exported figure.}}
  \label{fig:depth-width-curves}
\end{figure}

\begin{figure}[H]
  \centering
  \pendingplot[0.18\textheight]{Depth $\times$ width heatmap of the endpoint
  statistic, and the same heatmap expressed as a difference from the $(2,128)$
  cell.}
  \caption{Screen A. Depth-by-width heatmap.
  \pending{replace the box with the exported figure.}}
  \label{fig:depth-width-heatmap}
\end{figure}

\begin{figure}[H]
  \centering
  \pendingplot[0.18\textheight]{Actor parameter count against survival, one
  marker per cell, so a cell that matches the reference at a fraction of the
  parameters is visible directly.}
  \caption{Screen A. Cost against survival.
  \pending{replace the box with the exported figure.}}
  \label{fig:depth-width-cost}
\end{figure}

\subsection*{Reading}

\pending{answer four questions explicitly.} First, whether the $K=1$ row
matches $K=2$: if it does, one hop of message passing is sufficient on a local
\texttt{bus14} graph and the second layer is paying for nothing. Second, whether
$K=3$ recovers, matches, or degrades relative to $K=2$, which bears directly on
the Screen C finding that the hierarchy augmentations fail with only two layers
(Section~\ref{sec:screen-bus-structure}). Third, whether there is a width
threshold: a monotone width effect and a cliff between two adjacent widths lead
to different conclusions about what the graph vector has to carry. Fourth,
whether the replication check in Table~\ref{tab:depth-width-replication} is
small relative to the differences claimed in Screens C and D; if it is not, the
margins in those screens need re-reading with the wider error bar.

\section{Screen B: Encoder Operator and Readout}
\label{sec:screen-encoder-pooling}

\emph{Runs in progress. The design below is final; the coverage, result, and
reading subsections are placeholders to be filled from the completed runs.}

\subsection*{Question}

Two independent choices have been held fixed at their defaults. The
message-passing operator can be GINE, which adds the projected edge vector to
the sender state inside every message, or GAT, which uses the edge vector to
weight how strongly each incoming message counts
(Appendix~\ref{app:gine-architecture}). The readout can be the mean over valid
nodes or a per-channel maximum
(Section~\ref{subsec:graph-mean-pooling}). The two interact in principle: a
softmax-weighted aggregation followed by a maximum readout selects twice, while
mean pooling after an unweighted sum preserves how many nodes are in a given
state. Which combination suits the busbar graph?

The busbar schema is the relevant setting for a second reason. Its edge vector
carries only the physical line block and no relation one-hot
(Table~\ref{tab:gine-gat-edge-inventory}), so GAT's attention here is computed
from line state --- loading, status, cooldown --- rather than from relation
identity. If attention helps anywhere, a graph whose edges differ only by
physical measurement is where it should show.

\subsection*{Design}

A $2\times2$ factorial over encoder $\in\{\text{GINE},\text{GAT}\}$ and readout
$\in\{\text{mean},\text{max}\}$ with three seeds per cell. Only three cells are
newly launched --- 9 runs in
\path{configs/gnn_graph_screening/stage3_encoder_pooling} --- because the
GINE/mean cell is the reference configuration of
Section~\ref{subsec:screening-reference-config} and is reused rather than
repeated. GAT uses four heads averaged rather than concatenated, so all four
cells return the same 128 features. Every other setting is inherited from the
reference configuration seed by seed.

Reusing a cell has a cost that must be carried into the reading: the GINE/mean
numbers come from a different launch than the other three, so any difference
against them contains whatever run-to-run drift exists between launches. That
drift is exactly what Screen A's replication check measures
(Table~\ref{tab:depth-width-replication}), so the two screens should be read
together.

\subsection*{Coverage}

\pending{completed steps per run and the common budget. GAT is more expensive
per step than GINE at equal width, so the two encoder rows may not reach the
same step count under a shared walltime limit; if they do not, compare at the
common budget and say so rather than at each run's own endpoint.}

\subsection*{Result}

\begin{table}[H]
  \centering
  \small
  \caption{Screen B. Test survival by encoder and readout, averaged over three
  seeds, with the seed-to-seed standard deviation in parentheses. The GINE/mean
  cell is inherited from the reference configuration. \pending{fill the three
  new cells using the same endpoint statistic as the inherited cell.}}
  \label{tab:encoder-pooling-grid}
  \setlength{\tabcolsep}{6pt}
  \begin{tabular}{lcc}
    \toprule
    & \multicolumn{2}{c}{Readout} \\
    \cmidrule(lr){2-3}
    Encoder & Mean & Max \\
    \midrule
    GINE & 97.31 (1.22)\textsuperscript{\dag} & -- \\
    GAT & -- & -- \\
    \bottomrule
  \end{tabular}
  \\[0.4em]
  {\small \textsuperscript{\dag}Reference configuration, full-test survival of
  the best checkpoint; not a new run.}
\end{table}

\begin{figure}[H]
  \centering
  \pendingplot{Aggregated test episodic-survival curves for the four
  encoder/readout combinations, one standard deviation band across the three
  seeds, with the inherited GINE/mean cell drawn in grey.}
  \caption{Screen B. Survival curves by encoder and readout.
  \pending{replace the box with the exported figure.}}
  \label{fig:encoder-pooling-curves}
\end{figure}

\begin{figure}[H]
  \centering
  \pendingplot[0.18\textheight]{Two panels: readout compared within each
  encoder, and encoder compared within each readout. If the same readout wins
  in both encoder panels, the two factors are acting independently.}
  \caption{Screen B. Factors compared within each other.
  \pending{replace the box with the exported figure.}}
  \label{fig:encoder-pooling-facets}
\end{figure}

\subsection*{Reading}

\pending{report the four cells, then answer three questions.} Whether either
factor has a main effect large relative to the seed spread; whether they
interact, which the $2\times2$ can show but not explain; and whether any gain
from GAT is worth its cost per step, since a policy that reaches the same
survival with a cheaper operator is the better default for the larger grids.
State explicitly whether the inherited GINE/mean cell was re-run for
confirmation, and if it was not, bound the comparison with the replication gap
from Screen A.

\section{Screen C: Structural Augmentations of the Busbar Graph}
\label{sec:screen-bus-structure}

\subsection*{Question}

Three optional augmentations add computational structure to the busbar graph
without adding any physical measurement: same-substation edges ($e$), which let
the busbars of one substation exchange messages directly; substation-summary
nodes ($n$), which insert one learned hub per substation; and virtual-node
readout ($v$), which replaces mean pooling with a learned graph-summary node.
Each is motivated by the same intuition --- shorten the path between related
busbars --- so the question is whether any of them earns its cost, and whether
they compose.

\subsection*{Design}

A complete $2^3$ factorial with three seeds per cell, 24 runs, in
\path{configs/gnn_graph_screening/stage2_structure}. The graph type is
\texttt{bus} throughout and the baseline cell is the unaugmented
\texttt{e0n0v0}, which is the reference configuration of
Section~\ref{subsec:screening-reference-config}. Cells are named by their three
switches, so \texttt{e1n0v1} means same-substation edges on, summary nodes off,
virtual-node readout on. Because the design is a cube rather than a grid, the
eight cells are reported in full before any marginal.

\subsection*{Coverage}

All 24 runs finished, each reaching at least 99.2\% of the 15M-step budget.
Endpoint numbers come from the full-test statistic, so each cell aggregates
three runs $\times$ 201 chronics. The saved best checkpoint is late in training
for most runs, between 13.7M and 14.9M steps, but not for all: the
\texttt{e0n1v0} seed-1 checkpoint was selected at 2.65M steps and the
\texttt{e1n1v1} seed-2 checkpoint at 9.87M, which is a symptom of those runs
never improving on an early peak.

\subsection*{Result}

\begin{figure}[H]
  \centering
  \includegraphics[width=\linewidth]{figures/gs_s2_structure_vs_baseline_curves.png}
  \caption{Screen C. Test episodic survival during training, one panel per
  augmented cell against the unaugmented \texttt{e0n0v0} baseline in grey. Thin
  lines are the three seeds, the thick line is their mean, and the band is one
  standard deviation.}
  \label{fig:structure-curves}
\end{figure}

\begin{figure}[H]
  \centering
  \includegraphics[width=0.86\linewidth]{figures/gs_s2_structure_cube.png}
  \caption{Screen C. The factorial as a cube: each vertex is one of the eight
  cells and each edge flips a single factor with the other two held fixed. The
  baseline \texttt{e0n0v0} is circled in red.}
  \label{fig:structure-cube}
\end{figure}

\begin{figure}[H]
  \centering
  \includegraphics[width=0.92\linewidth]{figures/gs_s2_structure_seed_heatmap.png}
  \caption{Screen C. The same eight cells as a heatmap of the full-test
  survival of each run's best checkpoint, with the three individual seed values
  printed under each cell mean.}
  \label{fig:structure-seed-heatmap}
\end{figure}

\begin{table}[H]
  \centering
  \small
  \caption{Screen C. Cells ranked by full-test survival, with the difference
  from the unaugmented \texttt{e0n0v0} baseline.}
  \label{tab:structure-ranking}
  \setlength{\tabcolsep}{5pt}
  \begin{tabular}{lrrrrr}
    \toprule
    Cell & Mean (\%) & Std (pp) & Min (\%) & Max (\%) & $\Delta$ baseline (pp) \\
    \midrule
    \texttt{e0n0v0} & 97.31 & 1.22 & 96.55 & 98.72 & $0.00$ \\
    \texttt{e1n0v0} & 96.20 & 3.94 & 91.66 & 98.68 & $-1.11$ \\
    \texttt{e0n0v1} & 95.63 & 6.19 & 88.52 & 99.82 & $-1.68$ \\
    \texttt{e0n1v0} & 93.51 & 6.39 & 86.72 & 99.40 & $-3.80$ \\
    \texttt{e1n1v0} & 90.01 & 4.92 & 85.42 & 95.19 & $-7.30$ \\
    \texttt{e1n0v1} & 89.92 & 9.42 & 80.48 & 99.33 & $-7.39$ \\
    \texttt{e0n1v1} & 68.91 & 5.77 & 62.59 & 73.91 & $-28.40$ \\
    \texttt{e1n1v1} & 66.11 & 33.95 & 27.86 & 92.68 & $-31.20$ \\
    \bottomrule
  \end{tabular}
\end{table}

\begin{table}[H]
  \centering
  \small
  \caption{Screen C. Main effect of each factor (12 runs per level) and the four
  single-factor flips of the cube, each one enabling that factor with the other
  two held fixed. \texttt{Spread} is the range of those four flips.}
  \label{tab:structure-effects}
  \setlength{\tabcolsep}{4pt}
  \begin{tabular}{lrrrlrr}
    \toprule
    Factor & Off (\%) & On (\%) & Effect & Flips (pp) & Mean flip & Spread \\
     & & & (pp) & & (pp) & (pp) \\
    \midrule
    $e$ (substation edges) & 88.84 & 85.56 & $-3.28$ & $-1.1,\,-5.7,\,-3.5,\,-2.8$ & $-3.28$ & 4.60 \\
    $v$ (virtual readout) & 94.26 & 80.14 & $-14.11$ & $-1.7,\,-24.6,\,-6.3,\,-23.9$ & $-14.11$ & 22.92 \\
    $n$ (summary nodes) & 94.77 & 79.63 & $-15.13$ & $-3.8,\,-26.7,\,-6.2,\,-23.8$ & $-15.13$ & 22.92 \\
    \bottomrule
  \end{tabular}
\end{table}

\subsection*{Reading}

No augmentation earns its place, and the result is unambiguous: the
unaugmented \texttt{e0n0v0} baseline is the best cell at 97.31\%, and all
twelve single-factor flips in Table~\ref{tab:structure-effects} are negative.
Not one of the three mechanisms helps in any context. The baseline is also the
most stable cell in the screen, with a seed spread of 1.22 points against
3.9--33.9 for every augmented cell, so the augmentations cost reliability as
well as survival. The training curves in Figure~\ref{fig:structure-curves}
separate the cells well before the budget ends, so this is not an artefact of
where the runs were cut.

The two hierarchy factors dominate, at $-15.13$ points for summary nodes and
$-14.11$ for virtual-node readout, and they are not independent: each is mild
in the cell where the other is off ($n$ costs 3.80 points at \texttt{e0v0}; $v$
costs 1.68 at \texttt{e0n0}) and severe once the other is on, with both $n1v1$
cells landing between 66\% and 69\%. The upper face of the cube in
Figure~\ref{fig:structure-cube} shows this directly, and the spread column of
Table~\ref{tab:structure-effects} makes it explicit: for $n$ and $v$ the four
flips range over 23 percentage points, so the main effect alone is not a usable
summary of either factor.
\\
\\
The interpretation is that the two hierarchy mechanisms are redundant rather
than complementary. Both insert a learned aggregation node above the busbars,
and the virtual node attaches to the summary nodes when they exist
(Figure~\ref{fig:virtual-node-readout}), so enabling both stacks two learned
hubs in series between the busbars and the readout. With only two
message-passing layers, the readout of the \texttt{n1v1} cells is separated
from most physical nodes by more hops than the encoder has layers. Screen A
varies exactly that depth, so its $K=3$ row is the direct test of this
explanation.
\\
\\
Same-substation edges are the mildest of the three, at $-3.28$ points on
average and within a narrow band: its four flips range only from $-1.11$ to
$-5.71$, so unlike $n$ and $v$ it behaves like a small, consistent cost rather
than a context-dependent one. Its best cell \texttt{e1n0v0} sits 1.11 points
below the baseline, which is inside the seed spread of both cells, and the
per-seed values in Figure~\ref{fig:structure-seed-heatmap} show why nothing
stronger can be claimed: two of its seeds reach 98.7\% and 98.3\% while the
third reaches only 91.7\%, straddling the baseline's much tighter
96.6--98.7\%. The honest reading is that the edges neither help nor clearly
hurt on their own, and that they add candidate edges for no measured return.
\\
\\
The operational conclusion is to keep the plain busbar graph with mean readout
for subsequent stages, and to treat any future hierarchy proposal as needing a
deeper encoder rather than another augmentation on top of two layers.

\section{Screen D: Generator and Load Message Direction}
\label{sec:screen-message-direction}

\subsection*{Question}

In the disaggregated graph, each generator and load is a node attached to its
busbar. Section~\ref{sec:heterogeneous-graph-representation} leaves the
direction of that attachment open: the asset may send messages to the busbar,
receive them from it, or both. The two directions are not symmetric in effect,
because the mean readout of Section~\ref{subsec:graph-mean-pooling} pools
asset nodes as well as busbars. Does the choice matter, and do the generator
and load relations behave independently?

This is the one screen that leaves the busbar schema, so it has its own
baseline cell rather than the reference configuration of
Section~\ref{subsec:screening-reference-config}.

\subsection*{Design}

A complete $3\times3$ factorial over the generator and load directions with
three seeds per cell: 27 runs, held in the configuration folder
\path{configs/gnn_graph_screening/heterogenous_message_direction}.
The graph type is \texttt{heterogeneous} throughout; the line and summary
directions stay \texttt{bidirectional} and are inert for this schema. The
encoder matches the reference configuration --- GINE, two layers, 128
features, mean readout, no augmentations --- so the only screened factors are
the two attachment directions. The baseline cell is bidirectional on both
relations. Directions are abbreviated \texttt{bi} (bidirectional), \texttt{a2b}
(asset to busbar), and \texttt{b2a} (busbar to asset).

\subsection*{Coverage}

All 27 runs finished. Every run reached at least 96\% of the 15M-step budget
and most reached 14.97M steps; the common budget is 14.432M steps, set by
\texttt{gbi\_la2b\_s1}. Endpoint numbers come from the full-test statistic, so
each cell aggregates three runs $\times$ 201 chronics. The saved best
checkpoint is not always late in training: for one seed of the
\texttt{a2b}/\texttt{a2b} cell it was selected at 4.31M steps, which is itself
a symptom of that cell's instability.

\subsection*{Result}

\begin{figure}[H]
  \centering
  \includegraphics[width=\linewidth]{figures/gs_hmd_direction_vs_baseline_curves.png}
  \caption{Screen D. Test episodic survival during training, one panel per
  direction pair against the bidirectional baseline in grey. Thin lines are the
  three seeds, the thick line is their mean, and the band is one standard
  deviation.}
  \label{fig:direction-curves}
\end{figure}

\begin{figure}[H]
  \centering
  \includegraphics[width=0.96\linewidth]{figures/gs_hmd_direction_seed_heatmap.png}
  \caption{Screen D. Full-test survival of the best checkpoint by generator
  direction (rows) and load direction (columns), with the three individual seed
  values printed under each cell mean. The baseline cell is
  \texttt{bi}/\texttt{bi}.}
  \label{fig:direction-seed-heatmap}
\end{figure}

\begin{table}[H]
  \centering
  \small
  \caption{Screen D. Direction pairs ranked by full-test survival, with the
  difference from the \texttt{bi}/\texttt{bi} baseline.}
  \label{tab:direction-ranking}
  \setlength{\tabcolsep}{5pt}
  \begin{tabular}{lrrrrr}
    \toprule
    Generator / load & Mean (\%) & Std (pp) & Min (\%) & Max (\%) & $\Delta$ baseline (pp) \\
    \midrule
    \texttt{b2a} / \texttt{b2a} & 97.94 & 2.20 & 95.40 & 99.25 & $+5.30$ \\
    \texttt{b2a} / \texttt{a2b} & 97.43 & 1.42 & 96.32 & 99.03 & $+4.79$ \\
    \texttt{a2b} / \texttt{b2a} & 94.47 & 7.53 & 85.78 & 99.03 & $+1.83$ \\
    \texttt{bi} / \texttt{bi} & 92.65 & 8.24 & 83.21 & 98.40 & $0.00$ \\
    \texttt{bi} / \texttt{a2b} & 91.35 & 5.54 & 86.52 & 97.40 & $-1.30$ \\
    \texttt{bi} / \texttt{b2a} & 87.76 & 11.04 & 77.64 & 99.53 & $-4.89$ \\
    \texttt{b2a} / \texttt{bi} & 84.93 & 13.01 & 70.11 & 94.46 & $-7.72$ \\
    \texttt{a2b} / \texttt{bi} & 83.16 & 12.75 & 72.27 & 97.19 & $-9.49$ \\
    \texttt{a2b} / \texttt{a2b} & 56.41 & 22.12 & 30.88 & 69.60 & $-36.24$ \\
    \bottomrule
  \end{tabular}
\end{table}

\subsection*{Reading}

Busbar-to-asset is the better direction for both relations. Averaging over the
other relation, the generator marginal is 93.43\% for \texttt{b2a}, 90.58\% for
\texttt{bi}, and 78.01\% for \texttt{a2b}; the load marginal follows the same
order at 93.39\%, 86.91\%, and 81.73\%. The strongest part of the result is not
the mean but the spread: the two best cells, \texttt{b2a}/\texttt{b2a} and
\texttt{b2a}/\texttt{a2b}, have seed standard deviations of 2.20 and 1.42
percentage points against 8.24 for the bidirectional baseline, and their worst
seeds (95.40\% and 96.32\%) are above the baseline's mean. The seed values
printed in Figure~\ref{fig:direction-seed-heatmap} show this at a glance: the
bottom row holds three tightly grouped seeds per cell, while the baseline cell
spans 83.2\% to 98.4\%. A gain of about five points that comes with a fourfold
reduction in seed spread is a different kind of finding from a gain of five
points at unchanged variance.

The failure case is equally clear. Making both relations asset-to-busbar
collapses the policy to 56.41\% with a 22.12-point spread and a worst seed at
30.88\%, which is 36 points below the baseline and the largest single effect
measured anywhere in this chapter. Its panel in
Figure~\ref{fig:direction-curves} separates from the baseline within the first
two million steps and never recovers, so this is a training failure rather than
a late collapse.

A plausible mechanism, consistent with the readout defined in
Section~\ref{subsec:graph-mean-pooling} but not directly measured here, is that
the mean pooling includes generator and load nodes. Under \texttt{b2a} an asset
node receives busbar context and its updated embedding enters the readout, so
the asset row carries information about its local neighborhood. Under
\texttt{a2b} the asset node never receives a message, so after $K$ layers its
embedding is still a projection of its own raw features, and it dilutes the
pooled vector with context-free rows; with both relations set that way, every
generator and load row is inert. Testing this would require re-running the
worst cell with controlled readout, which restricts pooling to controlled
nodes, and is not part of the present screen.

Two caveats bound the conclusion. It is measured on \texttt{bus14} only, and
the full-test statistic reports selected checkpoints, so the level of every
cell is optimistic even though the comparison between cells is paired.

\section{Campaign Status}
\label{sec:screening-status}

Table~\ref{tab:screening-campaign-status} lists every screening folder that
exists in the repository, so that configured work is not read as measured work.
Only Screens C and D carry results; Screens A and B are training.

\begin{table}[H]
  \centering
  \small
  \caption{State of the graph-screening campaign. ``Complete'' means a finished
  factorial with a seed-aggregated summary in this chapter.}
  \label{tab:screening-campaign-status}
  \setlength{\tabcolsep}{4pt}
  \begin{tabular}{L{0.34\textwidth}rL{0.44\textwidth}}
    \toprule
    Configuration folder & Runs & Status \\
    \midrule
    \path{stage3_depth_width} & 36 & Training: Screen A \\
    \path{stage3_encoder_pooling} & 9 & Training: Screen B \\
    \path{stage2_structure} & 24 & Complete: Screen C \\
    \path{heterogenous_message_direction} & 27 & Complete: Screen D \\
    \midrule
    \path{stage1_representation_normalization} & 12 & Run under an earlier explicit-line edge definition; superseded \\
    \path{stage1b_normalization_confirmation} & 9 & Superseded by the same change \\
    \path{stage1c_provisional_structure} & 7 & Provisional; replaced by Screen C \\
    \path{stage1d_heterogeneous_directions} & 8 & Provisional; replaced by Screen D \\
    \path{stage1e_heterogeneous_structure} & 8 & Not analysed \\
    \path{stabilize_training} & 6 & Configured, not analysed \\
    \path{stage3_encoder} & 0 & Placeholder, no configurations \\
    \path{stage4_confirmation} & 0 & Placeholder, no configurations \\
    \bottomrule
  \end{tabular}
\end{table}

The two completed screens fix two defaults for the remaining stages: the busbar
graph is used without structural augmentation, and any run on the disaggregated
graph attaches generators and loads with busbar-to-asset relations rather than
the bidirectional default. Both defaults are established on \texttt{bus14}
alone and would need re-checking before being carried to a larger grid. Screens
A and B may revise the encoder that both defaults were measured with, in which
case the two completed screens should be spot-checked at the revised encoder
rather than assumed to carry over.
